\section*{Hoofdstuk \ref{ch:g5:raw-materials-and-recycling} --- Van grondstof tot recycling}

\begin{solution}{exo:g5:raw-materials-and-recycling:1}
Natuurlijk: hout, wol, steen. Vervaardigd: glas, staal, plastic.
\end{solution}

\begin{solution}{exo:g5:raw-materials-and-recycling:2}
Papier: hout (bomen). Glas: zand (met soda en kalksteen). Katoen: de
katoenplant. Plastic: ruwe aardolie.
\end{solution}

\begin{solution}{exo:g5:raw-materials-and-recycling:3}
Een gesteente waaruit een metaal gewonnen wordt. Aluminium komt van
bauxiet.
\end{solution}

\begin{solution}{exo:g5:raw-materials-and-recycling:4}
De pot in de glasbak, de krant bij het papier en karton, het blikje bij
het metaal, de fles bij de plastic flessen.
\end{solution}

\begin{solution}{exo:g5:raw-materials-and-recycling:5}
Nieuwe \omterm{def:g5:raw-materials-and-recycling:raw-material}{grondstof}: het gerecyclede materiaal wordt gebruikt in plaats van
zand, \omterm{def:g5:raw-materials-and-recycling:raw-material}{erts} of olie uit de natuur.
\end{solution}

\begin{solution}{exo:g5:raw-materials-and-recycling:6}
Hernieuwbaar: hout (als er opnieuw wordt aangeplant), wol, katoen. Niet
hernieuwbaar: ruwe aardolie, ijzererts.
\end{solution}

\begin{solution}{exo:g5:raw-materials-and-recycling:7}
Glas maken uit zand is een \omterm{def:g5:irreversible-changes:chemical-change}{chemische verandering}: er ontstaat iets
nieuws, en er is geen eenvoudige weg terug.
\end{solution}

\begin{solution}{exo:g5:raw-materials-and-recycling:8}
Bijvoorbeeld: minderen --- fruit kopen zonder plastic verpakking;
hergebruiken --- een glazen pot opnieuw vullen; recyclen --- een blikje
in de metaalbak gooien.
\end{solution}

\begin{solution}{exo:g5:raw-materials-and-recycling:9}
Elk materiaal wordt op zijn eigen manier gerecycled: stukjes van een
ander materiaal zouden het gesmolten glas bederven.
\end{solution}

\begin{solution}{exo:g5:raw-materials-and-recycling:10}
Een twintigste van 20 eenheden: $20 \div 20 = 1$ eenheid energie,
ongeveer.
\end{solution}

\begin{solution}{exo:g5:raw-materials-and-recycling:11}
Recyclen kost nog altijd energie, machines en vervoer; een voorwerp dat
nooit gemaakt wordt, bespaart zijn \omterm{def:g5:raw-materials-and-recycling:raw-material}{grondstof} en al die energie.
\end{solution}
